
% SGA 3，第 VII 篇_B；英文重建本第一轮的中文译本。
% 范围： 推论 3.3.1 全文 依据文献局部页 71--72,
% 止于下节之前： 推论 3.3.2.
% 合并阅读本页： 584--585. 重版印刷页： 573--574.
% 范围内原版边码： 538.
% 范围内编者注： 129--130.
% 依据： Polo--Gille Exp7B-13oct24.pdf,
% SHA-256 D303690E0195654BBE7C3DAA5D89EA23E24CD709A3575217234AC19AE8836694.
% 参照了 Jacob Reinhold 的 CC BY 4.0 Markdown，作为对照来源。
% 本组件无图；所有数学内容均为原生 LaTeX。
\subsubsection*{3.3.1. 推论}
\phantomsection
\label{sga3:VIIB:3.3.1}
{\itshape 设 \(S\) 是 \(\mathbb Q\) 上的局部 Noether 概形，\(G\) 是平坦且局部有限型的 \(S\)-群概形。\footnote{%
\label{sga3:VIIB:nde:129}编者添加了原文遗漏的平坦性假设。}
若 \(\omega_{G/S}\) 是局部自由 \(\mathcal O_S\)-模，则 \(G\) 在单位截面的每一点处都在 \(S\) 上平滑。}
\textit{证明。}
设 \(x\) 是单位截面上的点，\(s\) 是其在 \(S\) 中的像，分别用 \(\mathcal O_x\) 和 \(\mathcal O_s\) 表示 \(x\) 和 \(s\) 处的局部环。\footnote{%
\label{sga3:VIIB:nde:130}以下论证细节由编者补充。}
由假设，\(\mathcal O_s\) 和 \(\mathcal O_x\) 都是 Noether 局部环。因此，二者的极大理想进拓扑与有限余长度理想所定义的“伪紧”拓扑一致。用
\[
  \widehat{\mathcal O}_x
  \qquad\text{与}\qquad
  \widehat{\mathcal O}_s
\]
表示它们关于这个拓扑的完备化。由 EGA~IV\textsubscript{4}，17.5.3，\(G\) 在 \(x\) 处于 \(S\) 上平滑，当且仅当 \(\widehat{\mathcal O}_x\) 在 \(\widehat{\mathcal O}_s\) 上形式平滑，这两个代数均赋予极大理想进拓扑。因此只须证明后一断言。
环 \(\widehat{\mathcal O}_x\) 与 \(\widehat{\mathcal O}_s\) 是 \SGARefLink{sga3:VIIB:1.2.6}{1.2.6} 中定义的形式簇 \(\widehat G/\widehat S\) 与 \(\widehat S\) 在 \(x\) 与 \(s\) 处的局部环。置
\[
  k=\widehat{\mathcal O}_s,
  \qquad
  H=\operatorname{Spf}\bigl(\widehat{\mathcal O}_x\bigr).
\]
则 \(H\) 是无穷小形式 \(k\)-簇。由于 \(\widehat G/\widehat S\) 的构造与积交换\textup{（同上）}，\(H\) 实际上是无穷小形式 \(k\)-群。
令 \(I\) 为 \(\mathcal O_x\) 的增广理想。由假设，
\[
  \omega_{G/S,x}=I/I^2
\]
是秩为 \(n\) 的有限秩自由 \(\mathcal O_s\)-模。由于 \(\mathcal O_s\) 是 Noether 环，模 \(\omega_{H/k}\) 是 \(\omega_{G/S,x}\) 的完备化，可与
\[
  \omega_{G/S,x}
  \otimes_{\mathcal O_s}
  \widehat{\mathcal O}_s.
\]
等同。因此它是秩为 \(n\) 的拓扑自由 \(k\)-模。由\SGARefLink{sga3:VIIB:3.3:cartier-theorem}{定理 3.3} 的蕴含\textup{(iv)}\(\Rightarrow\)\textup{\SGARefLink{sga3:VIIB:3.3:cartier-theorem:item:ii}{(ii)}}，
\[
  \widehat{\mathcal O}_x
  \simeq
  k[[\omega_{H/k}]]
  \simeq
  k[[t_1,\ldots,t_n]].
\]
最后，由 EGA~0\textsubscript{IV}，19.3.3，后一形式幂级数代数在 \(k\) 上形式平滑。推论得证。
这样，我们重新得到了对任意概形 \(S\) 上局部有限呈示的群概形独立获得的一个结果；参见\ VI\textsubscript{B}，\SGARefLink{prop:sga3-VIB-1.6}{1.6}。
